%% kapa: exam=1 round=33 year=2022 period=1 form=A subject=economics subject_ko=경제학원론 seq=14 number=54
%% source: Q-Net 감정평가사 기출문제 (cst00309) artlSeq=5209997 file=raw/1008/5209997/2022년도 제33회 감정평가사 1차시험 1교시(민법,경제학,부동산학).pdf
%% conversion: kordoc 4.24.51 -> tools/gen1.py (auto); review=proofread-sonnet; flags=none
\begin{question}{54}
두 재화 X, Y를 소비하는 갑의 효용함수가 U(X, Y) = X$^{0.3}$Y$^{0.7}$이다. 이에 관한 설명으로 옳지 않은 것은?
\begin{choices}
  \item 선호체계는 단조성을 만족한다.
  \item 무차별곡선은 원점에 대해 볼록하다.
  \item 효용을 극대화할 때, 소득소비곡선은 원점을 지나는 직선이다.
  \item 효용을 극대화할 때, 가격소비곡선은 X재 가격이 하락할 때 Y재의 축과 평행하다.
  \item 효용을 극대화할 때, 소득이 2배 증가하면 X재의 소비는 2배 증가한다.
\end{choices}
\end{question}
\answer{4}
